\section{What went against the release}
\label{sec:against}

The protocol wrote down seven conditions that would mean v0.7 is not an advancement, in advance,
because after the numbers are in every one of them has a reading that makes it sound like something
else. Two came out against the release. This section is where they are treated as what they are, and
it is placed before the corrections and the limitations rather than after them.

\subsection{The composite cell: v0.7 does not beat what phase~2 already had}
\label{sec:against-composite}

\frozen{} \textbf{does not beat} \composite{}: \textbf{\vsevenBtwoComposite{}~points
$[\vsevenBtwoCompositeLo, \vsevenBtwoCompositeHi]$ over \vsevenBtwoCompositeN{} games.} Both banks are
positive, so the sign replicates, but the interval contains zero.

The protocol states the consequence in advance and it is quoted rather than paraphrased: what an
earlier phase called ``\vsevenKthreeOverComposite{}
$[\vsevenKthreeOverCompositeLo, \vsevenKthreeOverCompositeHi]$ on one bank'' did not replicate, and on
this cell the release reduces to \composite{} --- ``which is a phase-2 result, not a v0.7 one. The
report must then say that the v0.7 contribution is the \emph{measurement programme} and the four
closed ledger entries, not the policy.''

That is what this paper says.

The arithmetic makes the cell interpretable rather than merely disappointing. \frozen{} minus
\composite{} is exactly the group \texttt{rtie=1} $+$ urgency-off $+$ \texttt{stall=12} added and
\texttt{m2=0} removed, and \S\ref{sec:results-attribution} measures \texttt{stall=12}'s contribution
as identically zero and \texttt{m2=0}'s leave-one-out drop as zero. The protocol states that group's
training value in advance as \vsevenTrainGroup{}
$[\vsevenTrainGroupLo, \vsevenTrainGroupHi]$. On holdout it is \vsevenBtwoComposite{}
$[\vsevenBtwoCompositeLo, \vsevenBtwoCompositeHi]$. \S\ref{sec:against-selection} is what the
selection-bias check makes of that pair, and the answer is not comfortable.

\paragraph{What this does and does not license.} It does not retract
\S\ref{sec:results-frontier}: the frozen configuration \emph{is} a certified advancement over the
\vsix{} frontier, and \fcheap{} and \composite{} are different objects --- the first is a point of the
\vsix{} frontier, the second is a configuration built during phase~2 of \emph{this} programme. What it
licenses is the narrower and more useful statement: \textbf{the v0.7 programme produced a real
advance over \vsix{}, and it produced it in phase~2. The architecture work of phase~3, the bake-off of
phase~4 and the freeze added no measurable strength on top of it.}

\subsection{The location test does not replicate}
\label{sec:against-location}

The protocol's third failure condition is that the lattice cannot locate the gain: if removing every
single component in turn costs less than a third of the whole, then the configuration has a number it
cannot explain, and the standing rule that ``v0.7 must not headline a gain it cannot attribute'' binds.

Under the replication rule of \S\ref{sec:methodology-intervals} this must be read per bank, and the two
banks disagree:

\begin{center}\small
\begin{tabular}{lrrlc}
\toprule
Bank & Whole & One third & Largest single drop & Verdict \\
\midrule
7090001 & \vsevenLocWholeOne{} & \vsevenLocThirdOne{} & \texttt{\vsevenLocAtOne{}} \vsevenLocBiggestOne{} & \vsevenLocVerdictOne{} \\
7090003 & \vsevenLocWholeThree{} & \vsevenLocThirdThree{} & \texttt{\vsevenLocAtThree{}} \vsevenLocBiggestThree{} & \vsevenLocVerdictThree{} \\
\bottomrule
\end{tabular}
\end{center}

The two banks' largest drops differ by \vsevenLocSplit{}~points, against a half-width on that
difference of comparable size: the split is real on the point estimates and is \emph{not itself
statistically resolved}. Under the rule that a claim whose sign does not replicate is reported as not
replicated whatever its pooled interval says, \textbf{the pooled verdict is NOT REPLICATED}
(\vsevenLocReplicates{}), and the honest reading is that this battery does not settle whether the gain
is located by the protocol's one-third rule.

The sub-additivity claim of \S\ref{sec:results-attribution} passes on both banks. It is the one-third
location test specifically that splits, and the distinction matters: the components compose
predictably, but which one carries the majority of the whole is bank-dependent at this power.

\subsection{The selection signature, and where it appears}
\label{sec:against-selection}

\vsevenSelK{} configurations were scored against the deployed policy before the freeze. At the lattice
cell size the per-cell standard deviation is \vsevenSelSigma{} points, so the expected maximum of
\vsevenSelK{} draws under the null that none of them differs from the incumbent is
$\sigma\sqrt{2\ln K} = \vsevenSelExpectedMax{}$ points.

\textbf{How much of the measured gain would be expected from selection alone: \vsevenSelOnHoldout{}
points.} The holdout banks were never available for selection --- not one of the \vsevenSelK{} choices
could have been made using any deal of them --- so the selection term on holdout is zero by
construction and not by measurement. That is the strongest structural guarantee in this paper and it
is worth stating plainly before the awkward part.

The awkward part is the check's other direction: whether the holdout estimate is systematically
smaller than the training estimate by about that amount.

\begin{table}[t]\centering
\caption{The four quantities the preregistered selection check tabulates, against their holdout
measurements. Three reproduce or exceed; one falls short. These are not every figure the protocol
states in advance --- the partner-regime predictions of \S\ref{sec:results-partners} are others, and
two of those also come in low.}
\label{tab:selection}
{\small\begin{tabular}{lrrr}
\toprule
Quantity & Training & Holdout & Shortfall \\
\midrule
B2.4 --- \texttt{rtie} + urgency-off + \texttt{stall}, as a group & $+0.78$ & $+0.15$ & $-0.63$ \\
B6 self-play row, opponent \texttt{v05} & $+2.94$ & $+4.49$ & $+1.55$ \\
B7 diagonal (self-play) & $+4.51$ & $+4.55$ & $+0.04$ \\
B7 off-diagonal (cross-play) & $+4.48$ & $+4.55$ & $+0.07$ \\
\bottomrule
\end{tabular}
}
\end{table}

\vsevenSelReproduce{} of the four reproduce or exceed their training values. The one that falls short
is the composite cell, by \vsevenSelWorstShortfall{} points against an expected-maximum-under-the-null
of \vsevenSelExpectedMax{} --- \textbf{and that cell is precisely the one measuring the group of keys
the freeze was selected over.} The shortfall is the largest of the four and it falls on the cell measuring what the freeze was
chosen over, at about the size the null predicts. It is not the only shortfall the protocol's
advance figures produce --- the one-seat upgrade against a \vsix{} opponent comes in at
\vsevenOneSeatSix{} against a predicted \vsevenTrainOneSeatSix{}, which \S\ref{sec:results-partners}
prints --- so the correct statement is that the largest shortfall is where selection happened, not
that it is the only one. The training interval for that group was
$[\vsevenTrainGroupLo, \vsevenTrainGroupHi]$; the holdout interval is
$[\vsevenBtwoCompositeLo, \vsevenBtwoCompositeHi]$.

This is not proof that the group's training value was entirely selection --- one quantity cannot carry
that --- but it is the pattern the check was written to look for, it is where the check said to look,
and it is consistent with \S\ref{sec:against-composite}'s reading rather than with an alternative one.

\subsection{The seven conditions, answered}
\label{sec:against-seven}

\begin{center}\small
\begin{tabular}{clp{0.52\linewidth}}
\toprule
\# & Condition & Result \\
\midrule
1 & the frontier cell & \textbf{not met.} \vsevenBtwoFcheap{} $[\vsevenBtwoFcheapLo, \vsevenBtwoFcheapHi]$, sign replicates \\
2 & the composite cell & \textbf{MET.} \vsevenBtwoComposite{} $[\vsevenBtwoCompositeLo, \vsevenBtwoCompositeHi]$ \\
3 & attribution: the location test & \textbf{DOES NOT REPLICATE.} \vsevenLocVerdictOne{} on one bank, \vsevenLocVerdictThree{} on the other \\
4 & brittleness under partner change or cross-play & not met. Minimum \vsevenSoneMin{}; cross-play gap \vsevenXpGap{} \\
5 & a fresh exploit & not met. Every arm loses; largest upper bound \vsevenBfourWorstUpper{} \\
6 & the gate & not met. All \vsevenGateRules{} rules pass \\
7 & the controls & not met. Every control behaves as specified \\
\bottomrule
\end{tabular}
\end{center}

\textbf{Two of the seven are met or unresolved: item~2 outright, and item~3 by non-replication.} The
other five are not met. A report that led with the five and buried the two would be describing a
different evaluation from the one that ran.
